\SourceLeaf{383}%
since they give us an indisputable measure of the accuracy of the operations themselves. Apart from errors in calculation, they can, to within imperceptibly small quantities, be nothing other than the combined errors made in the astronomical determinations of longitude, the determinations of azimuth and the measurements of angles in the triangulation network. These effects cannot, of course, be separated, but the existence of the total error independently of the irregularities in the Earth's figure is undeniable, even if this means substantially revising downwards whatever opinion one may have had of the \textit{absolute} accuracy attained in all three operations. Presumably each contributed its share. although we are inclined to attribute the greater half to the measured angles of the triangles. We do indeed find in this work most of the operations that enter into these comparisons. The angle measurements between Milan and Superga, however, which French engineers made in earlier years, are given only in abridged form and already adjusted, so that one cannot judge their accuracy at all. We do find in the \textit{Connaissance des tems}, 1829, p.\,288, that errors of up to \(6''8\) occur in the sums of the angles. The triangle formed by Superga and the old and new observatories in Turin (p.\,254), which is very important for transferring the orientation in connecting the Milan and Turin observatories, likewise appears to have been measured with less than the required accuracy. To clarify this very important matter, at least as far as the two observatories are concerned, it would be highly desirable for the astronomers there to establish a new geodetic connection between them. With two heliotropes and the use of the station at S.\,Bernardo di Fenera, visible from both observatories, this would be extremely easy. If even one of the other former intermediate stations can be seen from there, as there is little reason to doubt, the whole task would require the measurement of only four angles.

The eighth section gives us a very interesting and valuable piece of work, a complete repetition of the survey of an arc of latitude carried out by \textsc{Beccaria} in 1762\,.\,.\,.\,1764. As is well known, the result of this survey could never be brought into agreement with the degrees measured in other countries. The curvature of the arc between the endpoints Mondovi and Andrate
